\begin{frame}{Détail des formules - Norme sigma}

\textbf{Norme sigma :} Fonction de normalisation pour éviter les singularités

\begin{equation}
\|\mathbf{z}\|_\sigma = \frac{1}{\epsilon} \left( \sqrt{1 + \epsilon \|\mathbf{z}\|^2} - 1 \right)
\end{equation}

\textbf{Gradient sigma :} Version normalisée du vecteur
\begin{equation}
\sigma_\epsilon(\mathbf{z}) = \frac{\mathbf{z}}{1 + \epsilon \sigma_{\text{norm}}(\mathbf{z})}
\end{equation}

\textbf{Paramètres :}
\begin{itemize}
\item $\epsilon = 0.1$ : paramètre de régularisation
\item Évite les divisions par zéro quand $\|\mathbf{z}\| \rightarrow 0$
\end{itemize}

\end{frame}

\begin{frame}{Détail des formules - Fonction de connectivité}

\textbf{Fonction $\rho_h$ :} Assure la connectivité entre robots

\begin{equation}
\rho_h(s) = \begin{cases}
1 & \text{si } 0 \leq s \leq h \\
\frac{1}{2}\left(1 + \cos\left(\frac{\pi(s-h)}{1-h}\right)\right) & \text{si } h < s \leq 1 \\
0 & \text{sinon}
\end{cases}
\end{equation}

\textbf{Fonction $\sigma_1$ :} Saturation douce
\begin{equation}
\sigma_1(s) = \frac{s}{\sqrt{1 + s^2}}
\end{equation}

\textbf{Paramètres :}
\begin{itemize}
\item $h = 0.2$ : seuil de connectivité maximale
\item Transition douce vers zéro pour éviter les déconnexions brusques
\end{itemize}

\end{frame}

\begin{frame}{Détail des formules - Terme de formation $\mathbf{u}_i^{\alpha}$}

\textbf{Fonction potentielle $\phi_{\alpha}$ :} Force d'attraction/répulsion
\begin{equation}
\phi_{\alpha}(s,d) = \frac{1}{2} \rho_h\left(\frac{s}{\sigma_{\text{norm}}(c)}\right) \left[(a+b) \sigma_1(s - \sigma_{\text{norm}}(d) + e) + (a-b)\right]
\end{equation}

\textbf{Vecteur unitaire entre robots :}
\begin{equation}
\mathbf{n}_{ij} = \sigma_\epsilon(\mathbf{p}_j - \mathbf{p}_i)
\end{equation}

\textbf{Paramètres :}
\begin{itemize}
\item $a = b = 1$ : coefficients d'attraction/répulsion
\item $e = \frac{|a-b|}{\sqrt{4ab}} = 0$ : paramètre d'asymétrie
\item $c = 5$ : portée d'interaction entre robots
\item $d$ : distance désirée entre robots
\end{itemize}

\end{frame}

\begin{frame}{Détail des formules - Contrôleur PID}

\textbf{Terme de formation complet :}
\begin{equation}
\mathbf{u}_i^{\alpha} = \sum_{j \in \mathcal{N}_i} \left[ K_p \phi_{\alpha}(\|\mathbf{p}_j - \mathbf{p}_i\|_\sigma, d_{ij}) \mathbf{n}_{ij} + K_i \int_0^t \phi_{\alpha} \mathbf{n}_{ij} d\tau + K_d \frac{d}{dt}(\phi_{\alpha} \mathbf{n}_{ij}) \right]
\end{equation}

\textbf{Gains PID :}
\begin{itemize}
\item $K_p = 0.2$ : gain proportionnel (réaction immédiate à l'erreur)
\item $K_i = 0.05$ : gain intégral (élimination erreur statique)
\item $K_d = 0.0$ : gain dérivé (amortissement, désactivé)
\end{itemize}

\textbf{Rôle :} Maintien des distances désirées $d_{ij}$ entre robots voisins

\end{frame}

\begin{frame}{Détail des formules - Terme de navigation $\mathbf{u}_i^{\gamma}$}

\textbf{Guidage vers la cible :}
\begin{equation}
\mathbf{u}_i^{\gamma} = -c_{1,\gamma}(\mathbf{p}_i - \mathbf{p}_r)
\end{equation}

\textbf{Adaptation dynamique du gain :}
\begin{equation}
c_{1,\gamma} = f(\text{proximité}, \text{angles}, \text{mode})
\end{equation}

\textbf{Modes de fonctionnement :}
\begin{itemize}
\item \textbf{Mode normal :} $c_{1,\gamma} = 0.25$
\item \textbf{Mode rotation :} $c_{1,\gamma} = 0.2$
\item \textbf{Adaptation :} Réduction si formation non respectée
\end{itemize}

\textbf{Lissage temporel :} $\alpha = 0.8$ pour éviter les transitions brusques

\end{frame}

\begin{frame}{Détail des formules - Évitement d'obstacles}

\textbf{Terme d'évitement $\mathbf{u}_i^{\beta}$ :}
\begin{equation}
\mathbf{u}_i^{\beta} = c_{1,\beta} \sum_{k \in \mathcal{O}_i} \phi_{\beta}(\|\mathbf{p}_{ik} - \mathbf{p}_i\|_\sigma, d_{\beta}) \mathbf{n}_{ik}
\end{equation}

\textbf{Fonction répulsive :}
\begin{equation}
\phi_{\beta}(s, d_{\beta}) = \rho_h\left(\frac{s}{d_{\beta}}\right) \left(\sigma_1(s - d_{\beta}) - 1\right)
\end{equation}

\textbf{Facteur de pondération :}
\begin{equation}
b_{ik} = \rho_h\left(\frac{\|\mathbf{p}_{ik} - \mathbf{p}_i\|_\sigma}{d_{\beta}}\right)
\end{equation}

\textbf{Paramètres :}
\begin{itemize}
\item $c_{1,\beta} = 0.3$ : gain d'évitement
\item $d_{\beta}$ : distance de sécurité aux obstacles
\end{itemize}

\end{frame}

\begin{frame}{Détail des formules - Saturation des commandes}

\textbf{Fonction de saturation :}
\begin{equation}
\text{sat}(\mathbf{u}_i) = \begin{cases}
u_{i,\min} & \text{si } u_i < u_{i,\min} \\
u_i & \text{si } u_{i,\min} \leq u_i \leq u_{i,\max} \\
u_{i,\max} & \text{si } u_i > u_{i,\max}
\end{cases}
\end{equation}

\textbf{Limites physiques :}
\begin{itemize}
\item $u_{i,\min} = -0.04$ m/s : vitesse minimale
\item $u_{i,\max} = 0.04$ m/s : vitesse maximale
\end{itemize}

\textbf{Commande finale saturée :}
\begin{equation}
\mathbf{u}_i^{\text{final}} = \text{sat}(\mathbf{u}_i^{\alpha} + \mathbf{u}_i^{\gamma} + \mathbf{u}_i^{\beta})
\end{equation}

\end{frame}

\begin{frame}{Détail des formules - Adaptation de $c_1^\gamma$}

\textbf{Conditions critiques (arrêt quasi-total) :}
\begin{equation}
r_{\min} < 0.3 \text{ ou } r_{\max} > 2.2 \Rightarrow c_1^\gamma \approx 0
\end{equation}

\textbf{Conditions pour valeur nominale :}
\begin{align}
(1 - \tau_d) &\leq r_{\min} \text{ et } r_{\max} \leq (1 + \tau_d) \\
|\theta_{\text{erreur},j}| &\leq \tau_\theta \quad \forall j
\end{align}

\textbf{Réduction progressive :}
\begin{equation}
c_1^\gamma = c_{1,\text{base}}^\gamma \cdot \min(f_d, f_\theta)
\end{equation}

\textbf{Paramètres :} $\tau_d = \pm 15\%$, $\tau_\theta = \pm 15°$

\end{frame}

\begin{frame}{Détail des formules - Facteurs de réduction $c_1^\gamma$}

\textbf{Facteur distance :}
\begin{equation}
f_d = f_{\text{proximité}} \cdot f_{\text{séparation}} \cdot f_{d,\text{add}}
\end{equation}

\textbf{Facteur angle :}
\begin{equation}
f_\theta = f_{\text{erreur angle}} \cdot f_{\theta,\text{add}}
\end{equation}

\textbf{Réductions exponentielles :}
\begin{equation}
f_{\text{proximité}} = \begin{cases}
e^{-k_{\min} \cdot \delta_{\min}} & \text{si } r_{\min} < s_{\min} \\
1 & \text{sinon}
\end{cases}
\end{equation}

\textbf{Paramètres empiriques :} $s_{\min} = 0.90$, $s_{\max} = 1.10$, $k_{\min} = k_{\max} = 0.4$

\end{frame}
